\documentclass[10pt,a4paper]{amsart}
\usepackage{amsmath,amssymb,amsthm,graphicx}
% For commutative diagrams we use Xy-pic
\usepackage[all]{xy}

% We use 2cell for 2-commutative diagrams.
\xyoption{2cell}
% Theorem environments.
%
\theoremstyle{plain}
\newtheorem{theorem}[subsection]{Theorem}
\newtheorem{proposition}[subsection]{Proposition}
\newtheorem{lemma}[subsection]{Lemma}

\theoremstyle{definition}
\newtheorem{definition}[subsection]{Definition}
\newtheorem{example}[subsection]{Example}
\newtheorem{exercise}[subsection]{Exercise}
\newtheorem{situation}[subsection]{Situation}

\theoremstyle{remark}
\newtheorem{remark}[subsection]{Remark}
\newtheorem{remarks}[subsection]{Remarks}

\numberwithin{equation}{subsection}

% Macros
%
\def\lim{\mathop{\mathrm{lim}}\nolimits}
\def\colim{\mathop{\mathrm{colim}}\nolimits}
\def\Spec{\mathop{\mathrm{Spec}}}
\def\Hom{\mathop{\mathrm{Hom}}\nolimits}
\def\Ext{\mathop{\mathrm{Ext}}\nolimits}
\def\SheafHom{\mathop{\mathcal{H}\!\mathit{om}}\nolimits}
\def\SheafExt{\mathop{\mathcal{E}\!\mathit{xt}}\nolimits}
\def\Sch{\mathit{Sch}}
\def\Mor{\mathop{\mathrm{Mor}}\nolimits}
\def\Ob{\mathop{\mathrm{Ob}}\nolimits}
\def\Sh{\mathop{\mathit{Sh}}\nolimits}
\def\NL{\mathop{N\!L}\nolimits}
\def\CH{\mathop{\mathrm{CH}}\nolimits}
\def\proetale{{pro\text{-}\acute{e}tale}}
\def\etale{{\acute{e}tale}}
\def\QCoh{\mathit{QCoh}}
\def\Ker{\mathop{\mathrm{Ker}}}
\def\Im{\mathop{\mathrm{Im}}}
\def\Coker{\mathop{\mathrm{Coker}}}
\def\Coim{\mathop{\mathrm{Coim}}}

% Boxtimes
%
\DeclareMathSymbol{\boxtimes}{\mathbin}{AMSa}{"02}

%
% Macros for moduli stacks/spaces
%
\def\QCohstack{\mathcal{QC}\!\mathit{oh}}
\def\Cohstack{\mathcal{C}\!\mathit{oh}}
\def\Spacesstack{\mathcal{S}\!\mathit{paces}}
\def\Quotfunctor{\mathrm{Quot}}
\def\Hilbfunctor{\mathrm{Hilb}}
\def\Curvesstack{\mathcal{C}\!\mathit{urves}}
\def\Polarizedstack{\mathcal{P}\!\mathit{olarized}}
\def\Complexesstack{\mathcal{C}\!\mathit{omplexes}}
% \Pic is the operator that assigns to X its picard group, usage \Pic(X)
% \Picardstack_{X/B} denotes the Picard stack of X over B
% \Picardfunctor_{X/B} denotes the Picard functor of X over B
\def\Pic{\mathop{\mathrm{Pic}}\nolimits}
\def\Picardstack{\mathcal{P}\!\mathit{ic}}
\def\Picardfunctor{\mathrm{Pic}}
\def\Deformationcategory{\mathcal{D}\!\mathit{ef}}
\newenvironment{reference}{}{}
\newenvironment{slogan}{}{}
\newenvironment{history}{}{}
\usepackage{cleveref}
\crefname{originalsection}{original section}{original sections}
\Crefname{originalsection}{Original section}{Original sections}
\makeatletter
\newlabel{situation-regular-model}{{0C61}{source}{}{sourcepointer-1}{}}
\newlabel{situation-regular-model@cref}{{[originalsection][0][]0C61}{[1][0][]source}{}{}{}}
\makeatother
\crefname{originaldefinition}{original definition}{original definitions}
\Crefname{originaldefinition}{Original definition}{Original definitions}
\makeatletter
\newlabel{definition-top-genus}{{0C79}{source}{}{sourcepointer-2}{}}
\newlabel{definition-top-genus@cref}{{[originaldefinition][0][]0C79}{[1][0][]source}{}{}{}}
\makeatother
\crefname{originaldefinition}{original definition}{original definitions}
\Crefname{originaldefinition}{Original definition}{Original definitions}
\makeatletter
\newlabel{definition-numerical-type-model}{{0CA5}{source}{}{sourcepointer-3}{}}
\newlabel{definition-numerical-type-model@cref}{{[originaldefinition][0][]0CA5}{[1][0][]source}{}{}{}}
\makeatother
\crefname{originallemma}{original lemma}{original lemmas}
\Crefname{originallemma}{Original lemma}{Original lemmas}
\makeatletter
\newlabel{lemma-numerical-type-of-model}{{0CA4}{source}{}{sourcepointer-4}{}}
\newlabel{lemma-numerical-type-of-model@cref}{{[originallemma][0][]0CA4}{[1][0][]source}{}{}{}}
\makeatother
\crefname{originalsection}{original section}{original sections}
\Crefname{originalsection}{Original section}{Original sections}
\makeatletter
\newlabel{curves-section-genus-geometric-genus}{{0CE0}{source}{}{sourcepointer-5}{}}
\newlabel{curves-section-genus-geometric-genus@cref}{{[originalsection][0][]0CE0}{[1][0][]source}{}{}{}}
\makeatother
\crefname{originallemma}{original lemma}{original lemmas}
\Crefname{originallemma}{Original lemma}{Original lemmas}
\makeatletter
\newlabel{curves-lemma-bound-geometric-genus}{{0CE1}{source}{}{sourcepointer-6}{}}
\newlabel{curves-lemma-bound-geometric-genus@cref}{{[originallemma][0][]0CE1}{[1][0][]source}{}{}{}}
\makeatother
\crefname{originallemma}{original lemma}{original lemmas}
\Crefname{originallemma}{Original lemma}{Original lemmas}
\makeatletter
\newlabel{curves-lemma-genus-normalization}{{0CE4}{source}{}{sourcepointer-7}{}}
\newlabel{curves-lemma-genus-normalization@cref}{{[originallemma][0][]0CE4}{[1][0][]source}{}{}{}}
\makeatother
\crefname{originallemma}{original lemma}{original lemmas}
\Crefname{originallemma}{Original lemma}{Original lemmas}
\makeatletter
\newlabel{curves-lemma-bound-geometric-genus-curve}{{0CE5}{source}{}{sourcepointer-8}{}}
\newlabel{curves-lemma-bound-geometric-genus-curve@cref}{{[originallemma][0][]0CE5}{[1][0][]source}{}{}{}}
\makeatother
\crefname{originallemma}{original lemma}{original lemmas}
\Crefname{originallemma}{Original lemma}{Original lemmas}
\makeatletter
\newlabel{curves-lemma-genus-base-change}{{0BY9}{source}{}{sourcepointer-9}{}}
\newlabel{curves-lemma-genus-base-change@cref}{{[originallemma][0][]0BY9}{[1][0][]source}{}{}{}}
\makeatother
\crefname{originallemma}{original lemma}{original lemmas}
\Crefname{originallemma}{Original lemma}{Original lemmas}
\makeatletter
\newlabel{morphisms-lemma-scheme-theoretic-union}{{0C4J}{source}{}{sourcepointer-10}{}}
\newlabel{morphisms-lemma-scheme-theoretic-union@cref}{{[originallemma][0][]0C4J}{[1][0][]source}{}{}{}}
\makeatother
\crefname{originallemma}{original lemma}{original lemmas}
\Crefname{originallemma}{Original lemma}{Original lemmas}
\makeatletter
\newlabel{varieties-lemma-divisible}{{0C6P}{source}{}{sourcepointer-11}{}}
\newlabel{varieties-lemma-divisible@cref}{{[originallemma][0][]0C6P}{[1][0][]source}{}{}{}}
\makeatother
\crefname{originallemma}{original lemma}{original lemmas}
\Crefname{originallemma}{Original lemma}{Original lemmas}
\makeatletter
\newlabel{varieties-lemma-variety-with-smooth-rational-point}{{0CDW}{source}{}{sourcepointer-12}{}}
\newlabel{varieties-lemma-variety-with-smooth-rational-point@cref}{{[originallemma][0][]0CDW}{[1][0][]source}{}{}{}}
\makeatother
\title{The first cohomology of the reduction of the special fibre dominates the topological genus plus the geometric genus}
\author{The Stacks Project Authors}
\date{stacks-project:models.tex:local-snapshot-20261012; GNU FDL 1.2 or later}
\begin{document}
\maketitle
\noindent\textit{Editorial note.} The counted unit is the Stacks Project Lemma lemma-genus-reduction-bigger-than at source lines 5431-5446 together with its proof at source lines 5447-5538, a straight adjacent pairing in the chapter file, and it is the companion of the card on lemma-genus-reduction-smaller. The chapter is the upstream file models.tex; the pool directory carries the longer name stacks-models and is not used as the chapter name anywhere in this card.\par\medskip
\noindent\textit{Editorial note.} This card has no context span, and the situation that the statement names is declared as a pointer instead of being retained: the statement needs it only to fix notation, and the twelve references that resolution requires are declared below.\par\medskip
\noindent\textit{Editorial note.} The retained statement and proof name twelve results: four of this chapter, namely the situation, the definition of the topological genus, the definition of the numerical type of a model and the lemma relating a model to its numerical type; and eight of the Algebraic Curves, Varieties and Morphisms chapters. All twelve are declared as non-load-bearing navigation pointers with their stable Stacks tags looked up in the shared tag index, and none of their statements is reproduced inside the span. The situation is declared with the kind section, because the assembler's whitelist of kinds has no entry for a situation.\par\medskip
\noindent\textit{Editorial note.} No citation occurs in the retained spans, so no bibliography entry is supplied.\par\medskip
\noindent\textit{Allowed standard prerequisites (editorial).}\begin{itemize}
\item The regular-model situation, the topological genus of a numerical type and the geometric genus of a curve. The situation fixes a discrete valuation ring, a smooth projective curve over its fraction field, a regular proper model and the components of the special fibre with their multiplicities; the numerical type attached to the fibre carries a topological genus formed from the number of components and the positive intersection pairs. For a proper curve over a field the geometric genus is the dimension of the first cohomology of its structure sheaf, and it can also be read off from the normalisation of the curve together with the singularities of the reduction. The author's own definitions of the chapter and the results of the Algebraic Curves and Varieties chapters that the proof quotes are the allowed inputs.. The counted lemma is one half of the genus formula of the section, its companion giving the reverse inequality; together they identify the genus of the curve with the topological genus of the special fibre.
\end{itemize}\medskip
\ifdefined\hypertarget\hypertarget{sourcepointer-1}{}\fi
\noindent\textit{Source reference.} Original section 0C61: focused-seeds/AGP/sources/stacks-snapshot-20261011/tags.txt, source lines 15281--15281. Pointer to the situation that the retained statement names; the situation is not retained as a context span for this card, because the statement only needs it to fix notation, so its reference is declared as a pointer with the tag of the situation. The number is the stable Stacks tag of the target, looked up in the shared tag index.\par\medskip
\ifdefined\hypertarget\hypertarget{sourcepointer-2}{}\fi
\noindent\textit{Source reference.} Original definition 0C79: focused-seeds/AGP/sources/stacks-snapshot-20261011/tags.txt, source lines 15324--15324. Navigation pointer to the chapter's definition of the topological genus, named by the retained statement; the cited definition is an allowed input and its text is not reproduced here. The number is the stable Stacks tag of the target, looked up in the shared tag index.\par\medskip
\ifdefined\hypertarget\hypertarget{sourcepointer-3}{}\fi
\noindent\textit{Source reference.} Original definition 0CA5: focused-seeds/AGP/sources/stacks-snapshot-20261011/tags.txt, source lines 15425--15425. Navigation pointer to the chapter's definition of the numerical type associated to a model, named by the retained statement; the cited definition is an allowed input and its text is not reproduced here. The number is the stable Stacks tag of the target, looked up in the shared tag index.\par\medskip
\ifdefined\hypertarget\hypertarget{sourcepointer-4}{}\fi
\noindent\textit{Source reference.} Original lemma 0CA4: focused-seeds/AGP/sources/stacks-snapshot-20261011/tags.txt, source lines 15424--15424. Navigation pointer invoked by the retained proof for the relation between the model and its numerical type, in particular the data of the special fibre; the cited statement is an allowed input of the same section and its proof is not reproduced here. The number is the stable Stacks tag of the target, looked up in the shared tag index.\par\medskip
\ifdefined\hypertarget\hypertarget{sourcepointer-5}{}\fi
\noindent\textit{Source reference.} Original section 0CE0: focused-seeds/AGP/sources/stacks-snapshot-20261011/tags.txt, source lines 15560--15560. Allowed input of the Algebraic Curves chapter quoted by the retained statement for the notion of geometric genus; the pointer is registered so the reference resolves and the cited material is not reproduced inside the span. The number is the stable Stacks tag of the target, looked up in the shared tag index.\par\medskip
\ifdefined\hypertarget\hypertarget{sourcepointer-6}{}\fi
\noindent\textit{Source reference.} Original lemma 0CE1: focused-seeds/AGP/sources/stacks-snapshot-20261011/tags.txt, source lines 15561--15561. Allowed input of the Algebraic Curves chapter quoted by the retained proof for a bound on the geometric genus; the pointer is registered so the reference resolves and the cited statement is not reproduced inside the span. The number is the stable Stacks tag of the target, looked up in the shared tag index.\par\medskip
\ifdefined\hypertarget\hypertarget{sourcepointer-7}{}\fi
\noindent\textit{Source reference.} Original lemma 0CE4: focused-seeds/AGP/sources/stacks-snapshot-20261011/tags.txt, source lines 15564--15564. Allowed input of the Algebraic Curves chapter quoted by the retained proof for the genus of the normalisation of a curve; the pointer is registered so the reference resolves and the cited statement is not reproduced inside the span. The number is the stable Stacks tag of the target, looked up in the shared tag index.\par\medskip
\ifdefined\hypertarget\hypertarget{sourcepointer-8}{}\fi
\noindent\textit{Source reference.} Original lemma 0CE5: focused-seeds/AGP/sources/stacks-snapshot-20261011/tags.txt, source lines 15565--15565. Allowed input of the Algebraic Curves chapter quoted by the retained proof for the comparison between the geometric genus and the genus of a reduced curve; the pointer is registered so the reference resolves and the cited statement is not reproduced inside the span. The number is the stable Stacks tag of the target, looked up in the shared tag index.\par\medskip
\ifdefined\hypertarget\hypertarget{sourcepointer-9}{}\fi
\noindent\textit{Source reference.} Original lemma 0BY9: focused-seeds/AGP/sources/stacks-snapshot-20261011/tags.txt, source lines 15007--15007. Allowed input of the Algebraic Curves chapter quoted by the retained proof for the behaviour of the genus under base change of the field; the pointer is registered so the reference resolves and the cited statement is not reproduced inside the span. The number is the stable Stacks tag of the target, looked up in the shared tag index.\par\medskip
\ifdefined\hypertarget\hypertarget{sourcepointer-10}{}\fi
\noindent\textit{Source reference.} Original lemma 0C4J: focused-seeds/AGP/sources/stacks-snapshot-20261011/tags.txt, source lines 15229--15229. Allowed input of the Morphisms chapter quoted by the retained proof to write the reduction as the scheme-theoretic union of its components; the pointer is registered so the reference resolves and the cited statement is not reproduced inside the span. The number is the stable Stacks tag of the target, looked up in the shared tag index.\par\medskip
\ifdefined\hypertarget\hypertarget{sourcepointer-11}{}\fi
\noindent\textit{Source reference.} Original lemma 0C6P: focused-seeds/AGP/sources/stacks-snapshot-20261011/tags.txt, source lines 15304--15304. Allowed input of the Varieties chapter quoted by the retained proof for the divisibility of the degree of an invertible sheaf on a component by its residue field degree; the pointer is registered so the reference resolves and the cited statement is not reproduced inside the span. The number is the stable Stacks tag of the target, looked up in the shared tag index.\par\medskip
\ifdefined\hypertarget\hypertarget{sourcepointer-12}{}\fi
\noindent\textit{Source reference.} Original lemma 0CDW: focused-seeds/AGP/sources/stacks-snapshot-20261011/tags.txt, source lines 15556--15556. Allowed input of the Varieties chapter quoted by the retained proof for the consequences of a smooth rational point on a variety over a field, which is the hypothesis of the statement; the pointer is registered so the reference resolves and the cited statement is not reproduced inside the span. The number is the stable Stacks tag of the target, looked up in the shared tag index.\par\medskip
\begin{lemma}
\label{lemma-genus-reduction-bigger-than}
In Situation \ref{situation-regular-model} assume $X_k$ has a $k$-rational
point $x$ which is a smooth point of $X_k \to \Spec(k)$. Then
$$
\dim_k H^1((X_k)_{red}, \mathcal{O}_{(X_k)_{red}}) \geq
g_{top} + g_{geom}(X_k/k)
$$
where $g_{geom}$ is as in
Algebraic Curves, Section \ref{curves-section-genus-geometric-genus}
and $g_{top}$ is the topological genus
(Definition \ref{definition-top-genus})
of the numerical type associated to $X_k$
(Definition \ref{definition-numerical-type-model}).
\end{lemma}


\begin{proof}
We are going to prove the inequality
$$
\dim_k H^1(D, \mathcal{O}_D) \geq g_{top}(D) + g_{geom}(D/k)
$$
for all connected reduced effective Cartier divisors
$D \subset (X_k)_{red}$ containing $x$ by induction
on the number of irreducible components of $D$.
Here $g_{top}(D) = 1 - m + e$ where $m$ is the number of
irreducible components of $D$ and $e$ is the number of
unordered pairs of components of $D$ which meet.

\medskip\noindent
Base case: $D$ has one irreducible component. Then $D = C_i$
is the unique irreducible component containing $x$.
In this case $\dim_k H^1(D, \mathcal{O}_D) = g_i$
and $g_{top}(D) = 0$. Since $C_i$ has a $k$-rational smooth point
it is geometrically integral
(Varieties, Lemma \ref{varieties-lemma-variety-with-smooth-rational-point}).
It follows that $g_i$ is the genus of $C_{i, \overline{k}}$
(Algebraic Curves, Lemma \ref{curves-lemma-genus-base-change}).
It also follows that $g_{geom}(D/k)$ is the genus of the normalization
$C_{i, \overline{k}}^\nu$ of $C_{i, \overline{k}}$. Applying
Algebraic Curves, Lemma \ref{curves-lemma-genus-normalization}
to the normalization morphism $C_{i, \overline{k}}^\nu \to C_{i, \overline{k}}$
we get
\begin{equation}
\label{equation-genus-change-special-component}
\text{genus of }C_{i, \overline{k}} \geq
\text{genus of }C_{i, \overline{k}}^\nu
\end{equation}
Combining the above we conclude that
$\dim_k H^1(D, \mathcal{O}_D) \geq g_{top}(D) + g_{geom}(D/k)$
in this case.

\medskip\noindent
Induction step. Suppose we have $D$ with more than $1$ irreducible
component. Then we can write $D = C_i + D'$ where $x \in D'$ and
$D'$ is still connected. This is an exercise in graph theory we leave
to the reader (hint: let $C_i$ be the component of $D$ which is
farthest from $x$). We compute how the invariants change.
As $x \in D'$ we have $H^0(D, \mathcal{O}_D) = H^0(D', \mathcal{O}_{D'}) = k$.
Looking at the short exact sequence of sheaves
$$
0 \to \mathcal{O}_D \to \mathcal{O}_{C_i} \oplus \mathcal{O}_{D'}
\to \mathcal{O}_{C_i \cap D'} \to 0
$$
(Morphisms, Lemma \ref{morphisms-lemma-scheme-theoretic-union})
and using additivity of euler characteristics we find
\begin{align*}
\dim_k H^1(D, \mathcal{O}_D) - \dim_k H^1(D', \mathcal{O}_{D'})
& =
-\chi(\mathcal{O}_{C_i}) + \chi(\mathcal{O}_{C_i \cap D'}) \\
& =
w_i(g_i - 1) + \sum\nolimits_{C_j \subset D'} a_{ij}
\end{align*}
Here as in Lemma \ref{lemma-numerical-type-of-model} we set
$w_i = [\kappa_i : k]$, $\kappa_i = H^0(C_i, \mathcal{O}_{C_i})$,
$g_i$ is the genus of $C_i$, and $a_{ij} = (C_i \cdot C_j)$.
We have
$$
g_{top}(D) - g_{top}(D') = -1 +
\sum\nolimits_{C_j \subset D'\text{ meeting }C_i} 1
$$
We have
$$
g_{geom}(D/k) - g_{geom}(D'/k) = g_{geom}(C_i/k)
$$
by Algebraic Curves, Lemma \ref{curves-lemma-bound-geometric-genus}.
Combining these with our induction hypothesis, we
conclude that it suffices to show that
$$
w_i g_i - g_{geom}(C_i/k) +
\sum\nolimits_{C_j \subset D'\text{ meets } C_i} (a_{ij} - 1) - (w_i - 1)
$$
is nonnegative. In fact, we have
\begin{equation}
\label{equation-genus-change}
w_i g_i \geq [\kappa_i : k]_s g_i \geq g_{geom}(C_i/k)
\end{equation}
The second inequality by
Algebraic Curves, Lemma \ref{curves-lemma-bound-geometric-genus-curve}.
On the other hand, since $w_i$ divides $a_{ij}$
(Varieties, Lemma \ref{varieties-lemma-divisible})
it is clear that
\begin{equation}
\label{equation-change-intersections}
\sum\nolimits_{C_j \subset D'\text{ meets } C_i} (a_{ij} - 1) - (w_i - 1)
\geq 0
\end{equation}
because there is at least one $C_j \subset D'$ which meets $C_i$.
\end{proof}

\end{document}
